\scanpage{100}
\mspage{\textcircled{3}}
On sait que $R_T$ est défini à équivalence \inserted{près} (unique à isom.
unique près), et que pour tte catégorie cofibrée $\mathcal F$ sur
$(\mathrm{Cat})_{\lambda}$, on a une équivalence de catégories (unique à isom unique près)
\[
\underline{\mathrm{Hom}}_{\mathrm{fib}/(\mathrm{Cat}_\lambda)}(T_\lambda,\mathcal F)\simeq\mathcal F(R_T)\tag*{(4.3)}
\]
Faisant $\mathcal F=T'_\lambda$, on trouve
\[
\underline{\mathrm{Hom}}_{\deleted{\unsure{\mathrm{fib}\,\mathrm{Cat}}}\lambda}(T,T')\simeq T'(R_T)\quad\text{(unique à isom unique près)}\tag*{(4.4)}
\]
et faisant en particulier $T'=\tau$, on trouve
\[
\underbrace{\mathrm{Hom}_\lambda(T,\tau)}_{\text{défn }\ R_T^\lambda}\simeq
\underbrace{\tau(R_T)}_{R_T\ \text{défn}}\qquad\text{d'où}\tag*{\deleted{(4.5)}}
\]
\[
R_T=R_T^\lambda\quad\text{(catégorie des $\lambda$-Constructions)}\tag*{(4.5)}
\]
\inserted{On \unsure{s'écrira|écrira} donc indifféremment $R_T$ ou $R_T^\lambda$. Alors (4.1)
s'interprète (à équivalence près) comme (3.8).}
\deleted{\illegible{1} Faisant \unsure{maintenant} \unsure{dans}} Dans le cas envisagé ici, la
\deleted{\illegible{1} \unsure{d'ailleurs} \unsure{soit la}} cat. des $\lambda$-constructions est une $\mathcal U$-cat.

Faisant $C=(\mathrm{Ens})$ dans (4.1), on trouve
\[
T(\mathrm{Ens})\simeq\underline{\mathrm{Hom}}_\lambda(R_T^\lambda,(\mathrm{Ens}))
\deleted{\unsure{\simeq|\cong}(R_T^\lambda)^{\unsure{\circ|o}}}\tag*{(4.6)}
\]
[si $T$ n'est pas $\lambda$-représentable, on a seulement

\fbox{5} Il faut préciser quelles sont \inserted{(mod équivalence)} les catégories
\unsure{1|2}-cofibrées sur $\mathrm{Cat}_\lambda$ (resp. munies d'un hom dans $\unsure{\lambda|\tau}_\lambda^I$)
qui proviennent des \deleted{``théories de} ``types de structures \inserted{relatifs à $\lambda$}'' -- donc
\deleted{préciser} \inserted{expliciter} une définition de la cat. ($\lambda$-types) (1.5)
\inserted{On va d'abord \deleted{\unsure{poser}} un $I$ et regarder les $\lambda$-types sur $I$.}
On va préciser \deleted{\unsure{4|les}} ``règles de formation''

a) \inserted{(Structure vide)} On a, \deleted{de $\unsure{\tau^I}$} pour $I\in\mathcal U$,
$\tau^I\in\mathrm{Ob}\,(\lambda\text{-}\inserted{I\text{-}}\mathrm{type})$

b) (\uline{\unsure{Adjonction|Rajouter} des axiomes}) \inserted{(d'isomorphie)} Si $T$ est un ($\lambda$-type) \inserted{sur $I$}, et
on se donne une \inserted{petite} partie \inserted{$A$} de $\mathrm{Fl}(S_{T,I}^{\lambda})\ (\subset\mathrm{Fl}(R_T^\lambda))$, alors
% slanted
\margin{On introduit $S_{T,I}^{\lambda}\subset R_T^\lambda$ $\hookrightarrow$ sous catégorie pleine engendrée par $I$ \deleted{\unsure{par}} \unsure{lim|lieus} de type $\lambda$ si}
la sous-catégorie 2-cofibrée pleine de $\unsure{T_A|T|T'_A}$ \deleted{formée \unsure{des}} \inserted{telle} que
$T_A(C)$ $(\deleted{\unsure{\xi|\subset}}C\in\mathrm{Ob}\,\mathrm{Cat}_\lambda)$ soit formé des $\xi\in T(C)$ tels
que (3.3) transforme $A$ en des isom, est un $\lambda$-\inserted{$I$-}type

\margin{NB peut-être faut il aussi se donner la possibilité d'égaliser des couples de flèches....}
% archive sheet number: 99
